% ----------------------
\subsection{Section 8 (April 1829)}
% ----------------------
	\label{DC8}
	
	% -----
	\subsubsection{Historical Background}
	% -----
		\label{DC8-HB}
		
		% --
		\paragraph{JSP}
		% --
		
			``In April 1829, soon after JS and Oliver Cowdery met and began working together on the translation of the plates, Cowdery not only wanted to write but also `became exceedingly anxious to have the power to translate bestowed upon him.' Several experiences related to the translation may have intensified his desire, including a revelation JS dictated for Cowdery in early April. `If thou wilt inquire,' the revelation promised Cowdery, `thou shalt know mysteries which are great and marvelous,' and further: `Behold I grant unto you a gift if you desire of me, to translate even as my servant Joseph.'%
				\footnote{%
					% \label{}%
					See \hyperref[DC6-11]{DC 6:11} and \hyperref[DC6-25]{25}.
					}
			
			``As JS dictated the translation of the Book of Mormon to Cowdery, even the words that Cowdery recorded described the gift of translation. Soon after translation work began, JS dictated several passages describing other ancient records and the divine means of translating them. A king by the name of Limhi, for example, told a man named Ammon that he possessed `twenty-four plates...filled with engravings' that he could not decipher, nor did he know anyone who could. Ammon told Limhi that he knew of a man who could translate the plates: `for he hath wherewith that he can look, and translate all records that are of ancient date; and it is a gift from God. And the things are called interpreters....And whosoever is commanded to look in them, the same is called seer[,]...a revelator, and a prophet.'%
				\footnote{%
					% \label{}%
					See \hyperref[MOSIAH-8-9]{Mosiah 8:9}, \hyperref[MOSIAH-8-13]{13}, and \hyperref[MOSIAH-8-16]{16}.
					}
			
			``The revelation featured below assured Cowdery that he could translate if he asked `with an honest heart' and with faith, and it declared, `Behold I will tell you in your mind \& in your heart by the Holy Ghost.' By implication, the revelation indicated that the gift to translate was not unlike other spiritual gifts that he possessed. Cowdery's first gift, according to this text, was `the spirit of Revelation,' the same `spirit by which Moses brought the children of Israel through the red Sea on dry ground.' Cowdery's second gift was identified as `the gift of working with the sprout,' or rod. Like many of his contemporaries, Cowdery probably used divining rods to find water or minerals,%
				\footnote{%
					% \label{}%
					The JSP's note cites Bushman's \textit{Joseph Smith and the Beginnings of Mormonism}, 98.
					}
			and though this gift may have been a `thing of Nature,' the revelation confirmed it was also a gift from God.%
				\footnote{%
					% \label{}%
					The JSP's note here: ``This affirmation of Cowdery's use of a `rod' as a divine gift illustrates the compatibility some early Americans perceived between biblical religion and popular supernaturalism. `From the outset,' according to historian Robert Fuller, `Americans have had a persistent interest in religious ideas that fall well outside the parameters of Bible-centered theology...In order to meet their spiritual needs...[they] switched back and forth between magical and Christian beliefs without any sense of guilt or intellectual inconsistency.' (Fuller, \textit{Spiritual, but Not Religious}, 15, 17; see also Ashurst-McGee, `Pathway to Prophethood,' 126--148.''
					}
			
			``According to Revelation Book 1, JS dictated four revelations in April 1829, all of them associated with translation. While these texts have a closely related historical context, the precise order of their dictation is unknown. One of the four, a revelation that declared itself the \hyperref[DC7]{translation of an ancient Johannine parchment}, was arranged before this revelation in the Book of Commandments and in all editions of the Doctrine and Covenants. JS's history follows the order of the Doctrine and Covenants by suggesting that the translation of the parchment may have come before this second revelation to Cowdery. Nevertheless, John Whitmer's ordering in Revelation Book 1, the earliest extant compilation of revelations, places the revelation featured here before the parchment. In the absence of more definitive information, Whitmer's ordering is followed here [in the JSP].''%
				\footnote{%
					% \label{}%
					In other words, the order in the print copy of the JSP is: \hyperref[DC6]{DC 6}, \hyperref[DC10]{DC 10}, \hyperref[DC8]{DC 8}, \hyperref[DC7]{DC 7}, and \hyperref[DC9]{DC 9}, which closes out April 1829.
					}
		
	% -----
	\subsubsection{Versions}
	% -----
		\label{DC8-V}
	
		\begin{tabular}{p{1in}p{2in}p{1in}p{2in}}
			JSP		& \hyperref[EMS-1828-1829-JS-APR-1829-B]{April 1829}	& BC	& Chapter 7, 19--20 \\
			DC 1835 & Section 34, 161--162									& TS	& 3:18 (15 July 1842), 853--854 \\
			MS		& 3:8 (December 1842), 134--135							& DC 1844	&  Section 34, 237--238 \\
		\end{tabular}
		
		\medskip
		
		See the \href{https://drive.google.com/file/d/1goAvNz8jS4R8slYfMG_db55Ty2z_vmgK/view?usp=sharing}{sections comparison} document for more information about the version history.

	% -----
	\subsubsection{Section 8:1} % *DC8-1 
	% -----
		\label{DC8-1}

			\footnote{%
				% \label{}%
				The JSP version has this header: ``A Revelation to Oliver [Cowdery] he being desirous to know whether the Lord would grant him the gift of \sout{Revelation} \sout{\& th$\diamond$} Translation given in Harmony Susquehannah Pennsylvania.'' BC has, ``1 \textit{A Revelation given to Oliver [Cowdery], in Harmony, Pennsylvania, April, 1829.}'' DC 1835 has, ``\textit{Revelation given April, 1829.}'' DC 1844 has, ``\textit{Revelation given April, 1829.}''
				}%
		OLIVER Cowdery, verily, verily, I say unto you, that assuredly as the Lord liveth, who is your God and your Redeemer, even so surely shall you receive a knowledge of whatsoever things you shall ask in faith, with an honest heart, believing that you shall receive a knowledge concerning the engravings of old records, which are ancient, which contain those parts of my scripture of which has been spoken by the manifestation of my Spirit.

		\begin{framed}
		\scriptsize
			\begin{compactdesc}
				\item[JSP] 3, a knowledge concerning the engraveings of old Records which are ancient which contain those parts of my Scriptures of which hath been spoken by the manifestation of my Spirit
				\item[BC] OLIVER, verily, verily I say unto you, that assuredly as the Lord liveth, which is your God and your Redeemer, even so sure shall you receive a knowledge of whatsoever things you shall ask in faith, with an honest heart, believing that you shall receive a knowledge concerning the engravings of old records, which are ancient, which contain those parts of my scripture of which have been spoken, by the manifestation of my Spirit;
				\item[DC 1835] 1 Oliver Cowdery, verily, verily I say unto you, that assuredly as the Lord liveth, who is your God and your Redeemer, even so sure shall you receive a knowledge of whatsoever things you shall ask in faith, with an honest heart, believing that you shall receive a knowledge concerning the engravings of old records, which are ancient, which contain those parts of my scripture of which have been spoken, by the manifestation of my Spirit;
				\item[DC 1844] 1 Oliver Cowdery, verily, verily I say unto you, that assuredly as the Lord liveth, who is your God and your Redeemer, even so sure shall you receive a knowledge of whatsoever things you shall ask in faith, with an honest heart, believing that you shall receive a knowledge concerning the engravings of the old records, which are ancient, which contain those parts of my scripture of which have been spoken, by the manifestation of my Spirit: 
			\end{compactdesc}
		\end{framed}

	% -----
	\subsubsection{Section 8:2} % *DC8-2 
	% -----
		\label{DC8-2}

		Yea, behold, I will tell you in your mind and in your heart, by the Holy Ghost, which shall come upon you and which shall dwell in your \hyperref[HEART]{heart}.

		\begin{framed}
		\scriptsize
			\begin{compactdesc}
				\item[JSP] yea Behold I will tell you in your mind \& in your heart by the Holy Ghost which Shall come upon you \& which shall dwell in your heart
				\item[BC] yea, behold I will tell you in your mind and in your heart by the Holy Ghost, which shall come upon you and which shall dwell in your heart.
				\item[DC 1835] yea, behold I will tell you in your mind and in your heart by the Holy Ghost, which shall come upon you and which shall dwell in your heart
				\item[DC 1844] yea, behold I will tell you in your mind and in your heart by the Holy Ghost, which shall come upon you and which shall dwell in your heart.
			\end{compactdesc}
		\end{framed}

	% -----
	\subsubsection{Section 8:3} % *DC8-3 
	% -----
		\label{DC8-3}

		Now, behold, this is the spirit	of revelation;%
			\footnote{%
				% \label{}%
				Very interesting to think about the possible meanings of ``spirit'' in this phrase; see commentary on the phrase ``spirit of revelation'' \hyperref[SPIRIT-PHRASES]{here}.
				}
		behold, this is the spirit by which Moses brought the children of Israel through the Red Sea on dry ground.%
			\footnote{%
				% \label{}%
				Commentary from Jeffrey Holland, 2 March 1999, in the ``\href{https://speeches.byu.edu/talks/jeffrey-r-holland/cast-not-away-therefore-your-confidence/}{Cast Not Away Therefore Your Confidence}'' talk at BYU (more of the talk is relevant too, but this is a good selection): ``Face your doubts. Master your fears. ``Cast not away therefore your confidence.'' Stay the course and see the beauty of life unfold for you. To help us make our way through these experiences, these important junctures in our lives, let me draw from another scriptural reference to Moses. It was given in the early days of this dispensation when revelation was needed, when a true course was being set and had to be continued. Virtually everyone in the room knows the formula for revelation given in \hyperref[DC9]{section 9} of the Doctrine and Covenants---you know, the verses about studying it out in your mind and the Lord promising to confirm or deny. What most of us don't read in conjunction with this is the section that precedes it---section 8. In that revelation the Lord defined \textit{revelation}: \textit{I will tell you in your mind and in your heart, by the Holy Ghost, which shall come upon you and which shall dwell in your heart.} [I love the combination there of both mind and heart. God will teach us in a reasonable way and in a revelatory way---mind and heart combined, by the Holy Ghost.] \textit{Now, behold,} this is the spirit of revelation; behold, this is the spirit by which Moses brought the children of Israel through the Red Sea on dry ground. [D\&C 8:2--3; emphasis added] Question: Why would the Lord use the example of crossing the Red Sea as the classic example of ``the spirit of revelation''? Why didn't he use the First Vision? Or the example from the book of Moses we just used? Or the vision of the brother of Jared? Well, he could have used any of these, but he didn't. Here he had another purpose in mind. Usually we think of revelation as information. Just open the books to us, Lord, like: What was the political significance of the Louisiana Purchase or the essence of the second law of thermodynamics? It is obvious that when you see those kinds of questions on a test paper, you need revelation. Someone said prayer will never be eliminated from the schools so long as there are final examinations. But aside from the fact that you probably aren't going to get that kind of revelation---because in this Church we do not believe in ex nihilo creation, especially in exams---this is too narrow a concept of \textit{revelation}. May I suggest how section 8 broadens our understanding of section 9, particularly in light of these ``fights of affliction'' that Paul spoke of and that I have been discussing. First of all, revelation almost always comes in response to a question, usually an urgent question---not always, but usually. In that sense it does provide information, but it is urgently needed information, special information. Moses' challenge was how to get himself and the children of Israel out of this horrible predicament they were in. There were chariots behind them, sand dunes on every side, and just a lot of water immediately ahead. He needed information all right---what to do---but it wasn't a casual thing he was asking. In this case it was literally a matter of life and death. You will need information, too, but in matters of great consequence it is not likely to come unless you want it urgently, faithfully, humbly. Moroni calls it seeking ``with real intent'' (\hyperref[MORONI-10-4]{Moroni 10:4}). If you can seek that way, and stay in that mode, not much that the adversary can counter with will dissuade you from a righteous path. You can hang on, whatever the assault and affliction, because you have paid the price to---figuratively, at least---see the face of God and live.''
				}

		\begin{framed}
		\scriptsize
			\begin{compactdesc}
				\item[JSP] now Behold this is the spirit of Revelation Behold this is the spirit by which Moses brought the children of Israel through the red Sea on dry ground
				\item[BC] 2 Now, behold this is the Spirit of revelation:--- behold this is the Spirit by which Moses brought the children of Israel through the Red sea on dry ground:
				\item[DC 1835] 2 Now, behold this is the Spirit of Revelation: behold this is the Spirit by which Moses brought the children of Israel though the Red sea on dry ground:
				\item[DC 1844] 2 Now, behold this is the Spirit of Revelation: behold this is the Spirit by which Moses brought the children of Israel through the Red Sea on dry ground;
			\end{compactdesc}
		\end{framed}

	% -----
	\subsubsection{Section 8:4} % *DC8-4 
	% -----
		\label{DC8-4}

		Therefore this is thy gift; apply unto it, and blessed art thou, for it shall deliver you out of the hands of your enemies, when, if it were not so, they would slay you and bring your soul to destruction.

		\begin{framed}
		\scriptsize
			\begin{compactdesc}
				\item[JSP] therefore this is thy gift apply unto it \& blessed art thou for [it] shall deliver you out of the hands of your enemies when if it were not so they would sley thee \& bring thy soul to distruction
				\item[BC] therefore, this is thy gift; apply unto it and blessed art thou, for it shall deliver you out of the hands of your enemies, when, if it were not so, they would slay you and bring your soul to destruction.
				\item[DC 1835] therefore this is thy gift; apply unto it and blessed art thou, for it shall deliver you out of the hands of your enemies, when, if it were not so, they would slay you and bring your soul to destruction.
				\item[DC 1844] therefore this is thy gift: apply unto it and blessed art thou, for it shall deliver you out of the hands of your enemies, when, if it were not so, they would slay you and bring your soul to destruction.
			\end{compactdesc}
		\end{framed}

	% -----
	\subsubsection{Section 8:5} % *DC8-5 
	% -----
		\label{DC8-5}

		Oh, remember these words, and keep my commandments.  Remember, this is your gift.

		\begin{framed}
		\scriptsize
			\begin{compactdesc}
				\item[JSP] O remember these words \& keep my commandments remember this is thy gift
				\item[BC] 3 O remember, these words and keep my commandments. Remember this is your gift.
				\item[DC 1835] 3 O remember these words, and keep my commandments.--- Remember this is your gift.
				\item[DC 1844] 3 O remember these words, and keep my commandments. Remember this is your gift. 
			\end{compactdesc}
		\end{framed}

	% -----
	\subsubsection{Section 8:6} % *DC8-6 
	% -----
		\label{DC8-6}

		Now this is not all thy gift; for you have another gift, which is the
			\footnote{%
				% \label{}%
				See the JSP version, which mentions the ``sprout.'' From the 1828 dictionary for ``sprout,'' meanings \#1 and \#2: ``The shoot of a plant; a shoot from the seed or from the stump or from the root of a plant or tree''; ``A shoot from the end of a branch.'' Then from the same dictionary, meaning \#3: ``A young branch.'' Part of a JSP note here: ``In preparing the text of Revelation Book 1 for publication, Sidney Rigdon replaced `sprout' with `rod.' Green, flexible shoots or rods cut from hazel, peach, or cherry trees were sometimes used as divining rods.''
				}%
		gift of Aaron;%
		\footnote{%
			% \label{}%
			Clearly there are very important differences between this modern version and the original manuscript version in the JSP and the original published version in the BC. See below for a note about the word ``sprout'' in the JSP version; by the 1835 DC, the text of this verse had taken basically the form we know now. You can think about it at least two ways: that they were trying to sanitize the text, or remove some references to things that were part of frontier mystery culture or that were otherwise going to be unacceptable in a published version of the text; or, you can see it as an evolving understanding of Cowdery's role in relation to JS and the work as a whole, such that he has assumed the role of Aaron in certain ways in his life. I think the latter explanation is better and more interesting, and leads to more interesting possibilities; it's also a possible way of connecting to the Aaronic priesthood, and things like the ministering of angels (\hyperref[DC107-20]{DC 107:20}), which is an important aspect of the epistemology of revelation.
			
			To add to this, here is HDDC 185--189: ``There is one major variation in the text of this revelation as it appears in the Book of Commandments. Verse 3 (verses 6--8 in the current edition of the D\&C) is more explicit about the additional gift given to Oliver Cowdery. This verse reads as follows: [BC text ``O remember, these words...that you shall know.''] In the current text of this verse, the gift is called the gift of Aaron. One of the gifts given Aaron was that of performing miracles through the use of a staff or rod (See \hyperref[EXODUS-4-2]{Exodus 4:2--4}; \hyperref[EXODUS-7-9]{7:9--13, 19--21}; \hyperref[EXODUS-8-5]{8:5, 6, 16, 17}; \hyperref[EXODUS-9-23]{9:23}; \hyperref[EXODUS-10-13]{10:13}; \hyperref[NUMBERS-17-1]{Numbers 17:1--10} as examples). The rod mentioned in this revelation is not necessarily Aaron's rod, but one that can be used in a similar way. The Book of Mormon prophesied that a rod would be available to a modern Joseph of this day (\hyperref[2NEPHI-3-17]{2 Nephi 3:17}), and there is some evidence that a rod has been used ([reference to an Andrus commentary on the PGP]). Evidently, not only did Oliver Cowdery have access to it, but also Orson Hyde (TS [1 April 1842], 741). It may even be that there was more than one of these rods, as is testified to by Solmon F. Kimball: `Under the date of June 21, 1892, Sister Sarah M. Kimball signed her name to the following statements: ``At a Relief Society Meeting held April 28, 1842, I heard the Prophet Joseph make this statement. `While other leading men of the Church have been unrighteously aspiring, Heber C. Kimball has been true and is to me what John was to Jesus, my beloved disciple.' Bro. Kimball showed me a rod that the Lord through the Prophet Joseph had given to him. He said that when he wanted to find out anything that was his right to know, all he had to do was to knell down with the rod in his hand, and that sometimes the Lord would answer his questions before he had time to ask them.'' My mother, and my sister, Helen Mar, told me the same thing and added to it, that Pres. Young received a similar rod from the Lord at the same time. They claimed that these rods were given to them because they were the only ones of the original Twelve who had not lifted up their \sout{hearts} heels against the Prophet' (Statement of Solomon F. Kimball, undated, located in the HDC).''
			
			I haven't yet been able to find the original Solomon Kimball source. In any case, we now know much more about the use of ``rods'' anyway, so the source isn't that important.
			}
		behold, it has told you many things;

		\begin{framed}
		\scriptsize
			\begin{compactdesc}
				\item[JSP] now this is not all for thou hast another gift which is the gift of working with the sprout Behold it hath told you things
				\item[BC] Now this is not all, for you have another gift, which is the gift of working with the rod: behold it has told you things:
				\item[DC 1835] Now this is not all thy gift; for you have another gift, which is the gift of Aaron: behold it has told you many things:
				\item[DC 1844] Now this is not all thy gift; for you have another gift, which is the gift of Aaron: behold it has told you many things: 
			\end{compactdesc}
		\end{framed}

	% -----
	\subsubsection{Section 8:7} % *DC8-7 
	% -----
		\label{DC8-7}

		Behold, there is no other power, save the power of God, that can cause this gift of Aaron to be with you.

		\begin{framed}
		\scriptsize
			\begin{compactdesc}
				\item[JSP] Behold there is no other power save God that can cause this thing of Nature to work in your hands
				\item[BC] behold there is no other power save God, that can cause this rod of nature, to work in your hands,
				\item[DC 1835] behold there is no other power save the power of God that can cause this gift of Aaron to be with you;
				\item[DC 1844] behold there is no other power save the power of God that can cause this gift of Aaron to be with you; 
			\end{compactdesc}
		\end{framed}

	% -----
	\subsubsection{Section 8:8} % *DC8-8 
	% -----
		\label{DC8-8}

		Therefore, doubt not, for it is the gift of God; and you shall hold it in your hands, and do marvelous works; and no power shall be able to take it away out of your hands, for it is the work of God.

		\begin{framed}
		\scriptsize
			\begin{compactdesc}
				\item[JSP] for it is the work of God
				\item[BC] for it is the work of God; 
				\item[DC 1835] therefore, doubt not, for it is the gift of God, and you shall hold it in your hands, and do marvelous works; and no power shall be able to take it away out of your hands; for it is the work of God. 
				\item[DC 1844] therefore, doubt not, for it is the gift of God, and you shall hold it in your hands, and do marvelous works; and no power shall be able to take it away out of your hands; for it is the work of God. 
			\end{compactdesc}
		\end{framed}

	% -----
	\subsubsection{Section 8:9} % *DC8-9 
	% -----
		\label{DC8-9}

		And, therefore, whatsoever you shall ask me to tell you by that means, that will I grant unto you, and you shall have knowledge concerning it.

		\begin{framed}
		\scriptsize
			\begin{compactdesc}
				\item[JSP] \& therefore whatsoever ye shall ask to tell you by that means that will he grant unto you that ye shall know
				\item[BC] and therefore whatsoever you shall ask me to tell you by that means, that will I grant unto you, that you shall know.
				\item[DC 1835] And therefore, whatsoever you shall ask me to tell you by that means, that will I grant unto you and you shall have knowledge concerning it:
				\item[DC 1844] And therefore, whatsoever you shall ask me to tell you by that means, that will I grant unto you and you shall have knowledge concerning it: 
			\end{compactdesc}
		\end{framed}

	% -----
	\subsubsection{Section 8:10} % *DC8-10 
	% -----
		\label{DC8-10}

		Remember that without faith you can do nothing; therefore ask in faith.  Trifle not%
			\footnote{%
				% \label{}%
				For ``trifle,'' see \hyperref[DC32-5]{DC 32:5} and notes there.
				}
		with these things; do not ask for that which you ought not.

		\begin{framed}
		\scriptsize
			\begin{compactdesc}
				\item[JSP] remember that without faith ye can do nothing trifle not with these things do not ask for that which ye had not ought
				\item[BC] 4 Remember that without faith you can do nothing. Trifle not with these things. Do not ask for that which you ought not. 
				\item[DC 1835] remember, that without faith you can do nothing. Therefore, ask in faith. Trifle not with these things: do not ask for that which you ought not: 
				\item[DC 1844] remember, that without faith you can do nothing. Therefore, ask in faith. Trifle not with these things: do not ask for that which you ought not: 
			\end{compactdesc}
		\end{framed}

	% -----
	\subsubsection{Section 8:11} % *DC8-11 
	% -----
		\label{DC8-11}

		Ask that you may know the \hyperref[MYSTERY]{mysteries} of God, and that you may translate and receive knowledge from all those ancient records which have been hid up, that are sacred; and according to your faith shall it be done unto you.

		\begin{framed}
		\scriptsize
			\begin{compactdesc}
				\item[JSP] ask that ye may know the mysteries of God \& that ye may Translate all those ancient Records which have been hid up which are Sacred \& according to your faith shall it be done unto you
				\item[BC] Ask that you may know the mysteries of God, and that you may translate all those ancient records, which have been hid up, which are sacred, and according to your faith shall it be done unto you.
				\item[DC 1835] ask that you may know the mysteries of God, and that you may translate and receive knowledge from all those ancient records which have been hid up, that are sacred, and according to your faith shall it be done unto you. 
				\item[DC 1844] ask that you may know the mysteries of God, and that you may translate and receive knowledge from all those ancient records which have been hid up, that are sacred, and according to your faith shall it be done unto you. 
			\end{compactdesc}
		\end{framed}

	% -----
	\subsubsection{Section 8:12} % *DC8-12 
	% -----
		\label{DC8-12}

		Behold, it is I that have spoken it; and I am the same that spake unto you from the beginning.  Amen.

		\begin{framed}
		\scriptsize
			\begin{compactdesc}
				\item[JSP] Behold it is I that have spoken it \& I am the same which spake unto you from the begining amen
				\item[BC] 5 Behold it is I that have spoken it, and I am the same which spake unto you from the beginning:--- Amen.
				\item[DC 1835] Behold, it is I that have spoken it: and I am the same who spake unto you from the beginning. Amen.
				\item[DC 1844] Behold, it is I that have spoken it: and I am the same who spake unto you from the beginning: Amen.
			\end{compactdesc}
		\end{framed}